\providecommand{\sppepatchroot}{../data_code/cesm215/Sppe_patch}

\label{app:sppe-namelist-controls}

This document describes a set of revisions to the MG2 microphysics in CESM2.1.5, collectively called \textbf{\textit{SPPE}} (Simplified PPE). Most revisions are derived directly from recent literature, while others are supported by microphysical reasoning. With the exception of the substepping parameter, all can be regarded as simplifications. No parameters were refitted, and no new processes were added.

\subsection*{Hallett--Mossop secondary ice}

\paragraph*{\texttt{factor\_HM = 0}.}
This factor multiplies the Hallett--Mossop ice-splinter number source produced when snow rims cloud droplets in the active temperature window.

Recent controlled experiments failed to reproduce efficient Hallett--Mossop rime splintering and found no evidence of the hundreds of secondary ice particles per milligram of rime reported in early experiments, motivating its removal here \href{https://doi.org/10.5194/acp-24-5247-2024}{(Seidel et al., 2024)}.

\paragraph*{\texttt{onoff\_sip = .false.}.}
This logical switch prevents both branches of the temperature-dependent Hallett--Mossop calculation from running, so the secondary-ice number source is zero.

Setting both is redundant but ensures that the Hallett--Mossop source cannot execute in any case.

The relevant excerpt from \texttt{micro\_mg\_utils.F90.patch} is:

\begin{verbatim}
if ((t(i) < 270.16_r8) .and. (t(i) >= 268.16_r8) .and. &
    onoff_sip) then
   nsacwi(i) = 3.5e8_r8*(270.16_r8-t(i))/2.0_r8 &
        * psacws(i)*factor_HM
else if ((t(i) < 268.16_r8) .and. (t(i) >= 265.16_r8) .and. &
         onoff_sip) then
   nsacwi(i) = 3.5e8_r8*(t(i)-265.16_r8)/3.0_r8 &
        * psacws(i)*factor_HM
else
   nsacwi(i) = 0.0_r8
end if
\end{verbatim}

\subsection*{Freezing-point depression}

\paragraph*{\texttt{depr\_point\_frz = .false.}.}
CAM normally uses aerosol water activity (log) multiplied by ice supersaturation to decide whether solute freezing-point depression prevents immersion freezing and to modify the critical germ radius.

The patch initializes all three aerosol-pathway flags as active and applies that water-activity gate only when \texttt{depr\_point\_frz} is true; setting the parameter to false therefore bypasses the solubility gate rather than disabling immersion freezing.

Immersion freezing occurs after an aerosol particle has activated and become immersed in a cloud droplet, so the droplet water activity should be close to one and freezing-point depression should not suppress this pathway strongly.

Sensitivity tests with \textbf{\textit{SPPE}} showed that retaining the depression term artificially suppressed immersion freezing because CAM treats aerosols as internally mixed; the bias was especially apparent in regions dominated by highly hygroscopic aerosols, including the Southern Ocean. This produced a spurious north--south contrast in cloud glaciation.

The dust-a1 branch in \texttt{hetfrz\_classnuc.F90.patch} illustrates the behavior:

\begin{verbatim}
do_bc   = .true.
do_dst1 = .true.
do_dst3 = .true.

if (depr_point_frz) then
   if (aw(id_dst1)*supersatice > 1._r8) then
      rgimm_dust_a1 = 2*vwice*sigma_iw / &
           (kboltz*t*LOG(aw(id_dst1)*supersatice))
   else
      do_dst1 = .false.
   end if
end if
\end{verbatim}

\subsection*{Ice-number and dust-INP limits}

\paragraph*{\texttt{enable\_nimax = .false.}.}
This setting bypasses the legacy MG1/MG2 condition that clips a positive cloud-ice number tendency when the predicted ice number would exceed \texttt{nimax}.

The legacy limiter is an unphysical, ad hoc constraint that was disabled in newer corrected CESM configurations; the problem and the NoNimax replacement are discussed by \href{https://doi.org/10.1029/2025MS004967}{Duffy et al. (2025)} and \href{https://doi.org/10.1029/2021MS002776}{Zhu et al. (2022)}.

This setting does not include the later replacement cap of 100\,cm\textsuperscript{$-3$}; instead, \texttt{limfacdu} provides a process-level limit on the dust immersion-freezing source only.

\begin{verbatim}
if (enable_nimax .and. do_cldice .and. nitend(i,k) > 0._r8 &
    .and. ni(i,k) + nitend(i,k)*deltat > nimax(i,k)) then
   nitend(i,k) = max(0._r8, &
        (nimax(i,k)-ni(i,k))/deltat)
end if
\end{verbatim}

\paragraph*{\texttt{limfacdu = 0.05}.}
This dimensionless factor caps the dust immersion-freezing source at 5\% of the available dust number per call before CAM takes the minimum of that cap and the classical-nucleation-theory rate.

The choice follows the 5\% dust-INP limit described in Sect.~2.2.2 of \href{https://doi.org/10.5194/gmd-16-1735-2023}{Gettelman et al. (2023)} and replaces the broad ice-number limiter with a limit applied directly to the relevant ice-nucleation source.

One of the cloud-borne dust branches in \texttt{hetfrz\_classnuc.F90.patch} is:

\begin{verbatim}
if (do_dst1) frzduimm = frzduimm + MIN( &
     limfacdu*total_cloudborne_aer_num(id_dst1)/deltat, &
     total_cloudborne_aer_num(id_dst1)/deltat &
     *(1._r8-exp(-Jimm_dust_a1*deltat)))
\end{verbatim}

\subsection*{Numerical substepping and cirrus isolation}

\paragraph*{\texttt{micro\_mg\_num\_steps = 8}.}
This native CAM control advances MG microphysics eight times during each MG call and passes \texttt{dtime/8} to every substep rather than using one long microphysical step.

By default, the 1800\,s atmospheric interval is divided first into three CLUBB--macrophysics--microphysics steps and then into eight MG steps, yielding an effective MG step of 75\,s.

This is the Mg2Sub8 configuration for which \href{https://doi.org/10.1029/2021MS002776}{Zhu et al. (2022)} reported convergence of cloud-ice number and shortwave cloud feedback.

The control is native to \texttt{micro\_mg\_cam.F90}, not introduced by the curated patches:

\begin{verbatim}
num_steps = micro_mg_num_steps

do it = 1, num_steps
   call micro_mg_tend2_0( &
        mgncol, nlev, dtime/num_steps, &
        ... )
end do
\end{verbatim}

\paragraph*{\texttt{naai\_het\_also\_in\_mpc = .false.}.}
This switch prevents the aerosol-aware ice-number target \texttt{naai} from being imposed at mixed-phase temperatures; with false, that source is allowed only below 233.15\,K, approximately $-40\,^{\circ}$C.

The setting isolates the cirrus use of \texttt{naai} and does not disable the separate heterogeneous mixed-phase freezing tendencies.

The patched condition in \texttt{micro\_mg2\_0.F90.patch} is:

\begin{verbatim}
where (naai > 0._r8 .and. t < icenuct .and. &
       relhum*esl/esi > 1.05_r8 .and. &
       (naai_het_also_in_mpc .or. t < 233.15_r8))
   nnuccd = (naai-ni/icldm)/mtime*icldm
   nnuccd = max(nnuccd,0._r8)
end where
\end{verbatim}

\subsection*{Selected aerosol and freezing pathways}

The six \texttt{hclas} factors independently multiply black-carbon (\texttt{bc}) and dust (\texttt{du}) immersion (\texttt{imm}), contact (\texttt{cnt}), and deposition (\texttt{dep}) nucleation rates before the combined tendencies reach MG2.

The complete routing expression added by \texttt{hetfrz\_classnuc\_cam.F90.patch} is:

\begin{verbatim}
frzimm = hclas_bc_imm*frzbcimm + hclas_du_imm*frzduimm
frzcnt = hclas_bc_cnt*frzbccnt + hclas_du_cnt*frzducnt
frzdep = hclas_bc_dep*frzbcdep + hclas_du_dep*frzdudep
\end{verbatim}

\paragraph*{\texttt{hclas\_du\_dep = 0.0}.}
This value removes the dust deposition-nucleation contribution from \texttt{frzdep}, retaining immersion freezing as the physically relevant dust-INP pathway in mixed-phase clouds.

\paragraph*{\texttt{hclas\_du\_cnt = 0.0}.}
This value removes the dust contact-nucleation contribution from \texttt{frzcnt}, whose atmospheric importance remains poorly constrained.

The mixed-phase-cloud rationale for both dust settings follows \href{https://doi.org/10.1029/2021RG000745}{Burrows et al. (2022)}, who identify immersion freezing as the predominant primary-ice mode, deposition freezing as minor, and contact freezing as uncertain because internally mixed INPs generally activate into droplets rather than remaining interstitial.

\paragraph*{\texttt{hclas\_bc\_imm = 0.0}.}
This value removes the black-carbon immersion-freezing contribution from \texttt{frzimm}.

\paragraph*{\texttt{hclas\_bc\_dep = 0.0}.}
This value removes the black-carbon deposition-nucleation contribution from \texttt{frzdep}.

\paragraph*{\texttt{hclas\_bc\_cnt = 0.0}.}
This value removes the black-carbon contact-nucleation contribution from \texttt{frzcnt}.

All three black-carbon factors are zero because black carbon makes an unimportant contribution to global INP concentrations relevant to mixed-phase clouds, with only limited regional exceptions \href{https://doi.org/10.1073/pnas.2001674117}{(Schill et al., 2020)}.

Together with the separately specified \texttt{hclas\_du\_imm}, these five zeros isolate dust immersion freezing as the heterogeneous aerosol pathway retained and perturbed by \textbf{\textit{SPPE}}.

\paragraph*{\texttt{onoff\_contact\_freezing = .false.}.}
This switch disables CAM's legacy internal droplet--aerosol contact-freezing routine when MG is configured without the aerosol-aware \texttt{hetfrz\_classnuc} pathway.

Contact nucleation from \texttt{hetfrz\_classnuc} is controlled independently by \texttt{hclas\_du\_cnt} and \texttt{hclas\_bc\_cnt}, which are also zero here.

\paragraph*{\texttt{rain\_freeze = .false.}.}
This switch suppresses the Bigg-style heterogeneous freezing of rain below 269.15\,K, making its rain-to-snow mass and number tendencies zero.

Rain freezing is disabled because this source is not aware of simulated INPs and instead follows an imposed temperature dependence that can mask the physical temperature dependence of INP-mediated freezing.


The two logical guards added by \texttt{micro\_mg\_utils.F90.patch} are:

\begin{verbatim}
if (qcic(i) >= qsmall .and. t(i) < 269.15_r8 .and. &
    onoff_contact_freezing) then
   ...
end if

if (t(i) < 269.15_r8 .and. qric(i) >= qsmall .and. &
    rain_freeze) then
   nnuccr(i) = pi*nric(i)*bimm &
        *(exp(aimm*(tmelt-t(i)))-1._r8)/lamr(i)**3
else
   mnuccr(i) = 0._r8
   nnuccr(i) = 0._r8
end if
\end{verbatim}

\subsection*{Phase of detrained condensate}

\paragraph*{\texttt{detrainment\_ramp\_liq = .true.}.}
Between 238.15 and 268.15\,K, this setting replaces the standard linear liquid--ice detrainment ramp with \texttt{dum1=0}, so the newly detrained condensate in that interval is assigned entirely to liquid rather than partitioned between liquid and ice.

Convectively detrained ice is therefore turned off in the mixed-phase interval because its imposed temperature ramp is not aware of simulated INPs and would mask the physical temperature dependence of INP-mediated ice formation.

Temperatures below 238.15\,K still use \texttt{dum1=1} and therefore retain all-ice detrainment, while temperatures above 268.15\,K already use all-liquid detrainment.

The same switch is applied in \texttt{clubb\_intr.F90.patch}, \texttt{macrop\_driver.F90.patch}, and the MG2 detrainment logic so the CLUBB and non-CLUBB paths use the same phase choice.

The CLUBB patch shows both the switch and the resulting phase partition:

\begin{verbatim}
if (state1%t(i,k) > 268.15_r8) then
   dum1 = 0.0_r8
else if (state1%t(i,k) < 238.15_r8) then
   dum1 = 1.0_r8
else
   if (detrainment_ramp_liq) then
      dum1 = 0.0_r8
   else
      dum1 = (268.15_r8-state1%t(i,k))/30._r8
   end if
end if

ptend_loc%q(i,k,ixcldliq) = dlfzm(i,k) &
     + dlf2(i,k)*(1._r8-dum1)
ptend_loc%q(i,k,ixcldice) = difzm(i,k) + dlf2(i,k)*dum1
\end{verbatim}

\subsection*{Combined \textbf{\textit{SPPE}} interpretation}

Collectively, the settings retain dust immersion freezing with a 5\% dust-number cap, remove the legacy ice-number limiter, and use eight MG2 substeps for numerical convergence.

They also keep aerosol-aware cirrus nucleation out of mixed-phase temperatures and force detrained condensate to remain liquid throughout the mixed-phase ramp.

Secondary ice, legacy contact freezing, rain freezing, dust contact/deposition nucleation, and all three black-carbon nucleation modes are disabled in this reference configuration.

\subsection*{Decision summary}

Table~\ref{tab:sppe-namelist-decisions} summarizes the physical or numerical rationale for every fixed control discussed above.

\refstepcounter{table}\label{tab:sppe-namelist-decisions}
\begin{longtable}{@{}>{\raggedright\arraybackslash}p{0.27\textwidth}>{\raggedright\arraybackslash}p{0.47\textwidth}>{\raggedright\arraybackslash}p{0.18\textwidth}@{}}
\toprule
\textbf{Decision} & \textbf{Short rationale} & \textbf{Reference} \\
\midrule
\endfirsthead
\toprule
\textbf{Decision} & \textbf{Short rationale} & \textbf{Reference} \\
\midrule
\endhead
\texttt{factor\_HM = 0} & Disable an efficient HM source that recent controlled experiments failed to reproduce. & \href{https://doi.org/10.5194/acp-24-5247-2024}{Seidel et al. (2024)} \\
\texttt{onoff\_sip = .false.} & Redundantly ensure that the HM source cannot execute. & \href{https://doi.org/10.5194/acp-24-5247-2024}{Seidel et al. (2024)} \\
\texttt{depr\_point\_frz = .false.} & Avoid artificial suppression after aerosol activation, when droplet water activity should be close to one. & \textbf{process-isolation} \\
\texttt{enable\_nimax = .false.} & Remove the unphysical legacy ice-number limiter used in the released CESM2 microphysics. & \href{https://doi.org/10.1029/2025MS004967}{Duffy et al. (2025)}; \href{https://doi.org/10.1029/2021MS002776}{Zhu et al. (2022)} \\
\texttt{limfacdu = 0.05} & Replace the broad ice limiter with a 5\% dust-number cap on the immersion-freezing source. & \href{https://doi.org/10.5194/gmd-16-1735-2023}{Gettelman et al. (2023)} \\
\texttt{micro\_mg\_num\_steps = 8} & Use a converged microphysical substep length. & \href{https://doi.org/10.1029/2021MS002776}{Zhu et al. (2022)} \\
\texttt{naai\_het\_also\_in\_mpc = .false.} & Prevent the cirrus ice-number target from leaking into mixed-phase clouds. & \textbf{process-isolation} \\
\texttt{hclas\_du\_dep = 0.0} & Remove the minor deposition pathway from mixed-phase dust nucleation. & \href{https://doi.org/10.1029/2021RG000745}{Burrows et al. (2022)} \\
\texttt{hclas\_du\_cnt = 0.0} & Remove the poorly constrained dust contact-freezing pathway. & \href{https://doi.org/10.1029/2021RG000745}{Burrows et al. (2022)} \\
\texttt{hclas\_bc\_imm = 0.0} & Remove the globally unimportant black-carbon immersion source. & \href{https://doi.org/10.1073/pnas.2001674117}{Schill et al. (2020)} \\
\texttt{hclas\_bc\_dep = 0.0} & Remove the globally unimportant black-carbon deposition source. & \href{https://doi.org/10.1073/pnas.2001674117}{Schill et al. (2020)} \\
\texttt{hclas\_bc\_cnt = 0.0} & Remove the globally unimportant black-carbon contact source. & \href{https://doi.org/10.1073/pnas.2001674117}{Schill et al. (2020)} \\
\texttt{onoff\_contact\_freezing = .false.} & Disable the legacy contact routine because activated, internally mixed INPs are generally already immersed in droplets. & \href{https://doi.org/10.1029/2021RG000745}{Burrows et al. (2022)} \\
\texttt{rain\_freeze = .false.} & Remove INP-unaware rain freezing based on an imposed temperature dependence. & \textbf{process-isolation} \\
\texttt{detrainment\_ramp\_liq = .true.} & Remove INP-unaware detrained ice based on an imposed temperature ramp. & \textbf{process-isolation} \\
\bottomrule
\end{longtable}

\section*{Outlook}
\addcontentsline{toc}{section}{Outlook}

\subsection*{Droplet freezing timing in CNT}

In the classical nucleation theory (CNT) droplet-freezing scheme, the maximum dust immersion-freezing source depends explicitly on the model microphysics time step through \nolinkurl{total_cloudborne_aer_num(id_dst1)/deltat}.

Increasing the number of microphysical substeps decreases \texttt{deltat} and therefore artificially increases this source limit.

A physically based \texttt{deltat\_CNT} should instead represent the time required to replenish activated ice-nucleating particles after droplet freezing, analogous to the replenishment interval in a typical cloud-chamber CNT experiment.

This replenishment timescale could be constrained against cloud-phase metrics rather than being tied to the numerical microphysics time step.

Sensitivity tests indicate that \texttt{deltat\_CNT} should be approximately two hours.

Within the source-limited regime, changing \texttt{deltat\_CNT} is mathematically equivalent to changing the prescribed ice-nucleating-particle efficiency.
%
